% Prose rendering of PrimeTensor/Algebra.lean.
% Fragment: no \documentclass, no preamble, no document environment.

\section{Multiplicative algebra on the completed carrier}
\label{Algebra:sec}

\leanfile{PrimeTensor/Algebra.lean}

\subsection*{What this file establishes}

\texttt{PrimeTensor/Completion.lean} produced the carrier $\name{MulReal}$ as a
quotient of Cauchy streams and left it with no operations at all.  This file
supplies them.  Multiplication and inversion are lifted from $\name{MulRat}$
to streams and then through asymptotic equivalence, and the resulting
$\name{MulReal}$ is shown to be a commutative group: associative, commutative,
unital, and with $a \cdot a^{-1} = 1$ for every element.

The lift runs in three strata, and it is worth naming them because the work is
very unevenly distributed between them.

\begin{enumerate}[label=(\roman*)]
\item \emph{The operations respect the scale.}  Two theorems,
  $\name{scaleWithin\_mul}$ and $\name{scaleWithin\_inv}$.  This is the whole
  analytic content of the file.
\item \emph{Hence they act on streams and respect asymptotic equivalence.}
  Four routine consequences, all of them the same anchor-juggling already seen
  in the transitivity proof of \texttt{PrimeTensor/Completion.lean}.
\item \emph{Hence they descend to the quotient, and the group laws hold there.}
  Seven laws, none of which involves an anchor or a level at all: each holds
  \emph{termwise} on streams, and $\name{of\_pointwise}$ converts a termwise
  equality into an equality of classes.
\end{enumerate}

So the analysis is confined to stratum (i), two theorems, and the algebra of
stratum (iii) is in every case ``the same law one layer down''.  What makes
that separation possible is that the group laws are exact identities on
$\name{MulRat}$, so no error term survives to be estimated.

\subsection{Ratios of products and of inverses}

\begin{lemma}[\lean{MulRat.inv\_mul\_pair} and \lean{MulRat.inv\_one}]
\label{Algebra:lem:inv-mul-pair}
$(a \cdot b)^{-1} = a^{-1} \cdot b^{-1}$, and $1^{-1} = 1$.
\end{lemma}

\begin{leancode}
theorem inv_mul_pair (a b : MulRat) :
    (a * b)⁻¹ = a⁻¹ * b⁻¹ := by
  refine Quotient.inductionOn₂ a b ?_
  intro x y
  change
    Quotient.mk PrimeRatio.setoid ((x * y)⁻¹) =
      Quotient.mk PrimeRatio.setoid (x⁻¹ * y⁻¹)
  rfl
\end{leancode}

\begin{proof}
Both are $\ctor{rfl}$ once representatives are chosen.  On
$\name{PrimeRatio}$, multiplication is componentwise and inversion exchanges
the two components, so $(xy)^{-1}$ and $x^{-1}y^{-1}$ are the same pair
$\langle x.\name{lower}\cdot y.\name{lower},\;
 x.\name{upper}\cdot y.\name{upper}\rangle$ --- not merely equivalent modulo
the setoid, but syntactically identical.  Likewise $1 = \langle 1,1\rangle$ is
fixed by the exchange.
\end{proof}

\begin{lemma}[\lean{MulRat.mul\_four\_shuffle}, private]
\label{Algebra:lem:four-shuffle}
$(ab)(cd) = (ac)(bd)$.
\end{lemma}

\begin{proof}
Five rewrites, alternating associativity and the single commutation $bc = cb$.
As in \texttt{PrimeTensor/Scale/Laws.lean}, the regrouping is written out
because $\name{MulRat}$ carries $\ctor{Mul}$ and $\ctor{One}$ instances but no
$\ctor{CommMonoid}$, so $\ctor{ac\_rfl}$ is not available.
\end{proof}

\begin{lemma}[\lean{MulRat.ratio\_mul\_pair}]\label{Algebra:lem:ratio-mul}
$\name{ratio}(ac, bd) = \name{ratio}(a,b) \cdot \name{ratio}(c,d)$.
\end{lemma}

\begin{proof}
Unfolding $\name{ratio}$, the left side is $(ac)(bd)^{-1}$, which
Lemma~\ref{Algebra:lem:inv-mul-pair} turns into $(ac)(b^{-1}d^{-1})$, and
Lemma~\ref{Algebra:lem:four-shuffle} regroups into
$(ab^{-1})(cd^{-1})$.
\end{proof}

\begin{designnote}[The multiplicative product rule is exact]
\label{Algebra:note:exact-product}
Lemma~\ref{Algebra:lem:ratio-mul} is the statement that \emph{the relative
error of a product is the product of the relative errors}.  It is an identity,
with no remainder and no cross term.

Its additive counterpart is not.  Estimating
$\lvert ac - bd\rvert$ requires splitting
$ac - bd = a(c-d) + (a-b)d$ and then \emph{bounding $a$ and $d$}, which is why
the classical construction of $\mathbb{R}$ from Cauchy sequences has to prove
first that a Cauchy sequence is bounded, and why the proof that a product of
Cauchy sequences is Cauchy cites that lemma.  Nothing of the sort appears in
this file, or anywhere in this development: there is no boundedness lemma
because there is no cross term to bound.  The multiplicative presentation buys
this by measuring error as a ratio rather than a difference, and a ratio is
multiplicative on the nose.
\end{designnote}

\begin{lemma}[\lean{MulRat.ratio\_inv\_pair}]\label{Algebra:lem:ratio-inv}
$\name{ratio}(a^{-1}, b^{-1}) = \name{ratio}(b, a)$.
\end{lemma}

\begin{proof}
$\name{ratio}(a^{-1},b^{-1}) = a^{-1}(b^{-1})^{-1} = a^{-1}b = ba^{-1} =
\name{ratio}(b,a)$, using $\name{inv\_inv}$ and commutativity from
\texttt{PrimeTensor/MulRat.lean}.
\end{proof}

\subsection{The two analytic theorems}

\begin{theorem}[\lean{MulRat.scaleWithin\_mul}]\label{Algebra:thm:scalewithin-mul}
If
\[
  \name{ScaleWithin}(\ctor{succ}\,\ell,\, a,\, b)
  \quad\text{and}\quad
  \name{ScaleWithin}(\ctor{succ}\,\ell,\, c,\, d),
\]
then $\name{ScaleWithin}(\ell,\, a \cdot c,\, b \cdot d)$.
\end{theorem}

\begin{proof}
Identical in shape to $\name{scaleWithin\_comp}$ of
\texttt{PrimeTensor/Scale/Laws.lean}, with
Lemma~\ref{Algebra:lem:ratio-mul} in place of the composition of ratios.
Unfolding $\name{scalePow}$ at $\ctor{succ}\,\ell$ turns the four hypotheses
into squares below $\name{two}$.  For the forward conjunct, rewrite
$\name{ratio}(ac,bd) = \name{ratio}(a,b)\name{ratio}(c,d)$, distribute the
ladder over that product by $\name{scalePow\_mul}$, and apply
$\name{mul\_lt\_two\_of\_sq\_lt\_two}$.  For the backward conjunct, rewrite
instead $\name{ratio}(bd,ac) = \name{ratio}(b,a)\name{ratio}(d,c)$ and use the
two backward hypotheses.
\end{proof}

\begin{theorem}[\lean{MulRat.scaleWithin\_inv}]\label{Algebra:thm:scalewithin-inv}
$\name{ScaleWithin}(\ell, a, b)$ implies
$\name{ScaleWithin}(\ell, a^{-1}, b^{-1})$ --- at the \emph{same} level.
\end{theorem}

\begin{leancode}
theorem scaleWithin_inv {level : Depth} {a b : MulRat}
    (h : ScaleWithin level a b) :
    ScaleWithin level a⁻¹ b⁻¹ := by
  unfold ScaleWithin at h ⊢
  rw [ratio_inv_pair a b, ratio_inv_pair b a]
  exact ⟨h.2, h.1⟩
\end{leancode}

\begin{proof}
By Lemma~\ref{Algebra:lem:ratio-inv} the two oriented ratios of the inverted
pair are the two oriented ratios of the original pair, exchanged.  So the goal
is the hypothesis with its conjuncts swapped.
\end{proof}

\begin{designnote}[Inversion is an isometry, and why that is free]
\label{Algebra:note:isometry}
Theorem~\ref{Algebra:thm:scalewithin-inv} costs no level: closeness at $\ell$
goes to closeness at $\ell$, and the proof is a swap.  Multiplicative distance
is literally invariant under $x \mapsto x^{-1}$, which merely reverses
orientation, and $\name{ScaleWithin}$ was defined symmetrically in the two
orientations precisely so that the reversal is invisible.

Contrast the classical construction again.  There, showing that the reciprocal
of a Cauchy sequence is Cauchy is the delicate step: $x \mapsto 1/x$ is not
uniformly continuous near $0$, the estimate
$\lvert 1/a - 1/b\rvert = \lvert a-b\rvert / \lvert ab\rvert$ blows up, and one
must first prove that a sequence not asymptotic to $0$ is eventually bounded
\emph{away} from $0$, then restrict to such sequences, then check the choice is
independent of representative.  None of that appears here, and the reason is
structural rather than clever: $\name{MulRat}$ is the strictly positive
rationals, so the singular point is not in the carrier to begin with.  The
completion is a group because the space being completed already was one, with
no punctured-domain bookkeeping to carry through the quotient.
\end{designnote}

\begin{remark}[One level for multiplication, none for inversion]
\label{Algebra:rem:level-cost}
The asymmetry is exactly right.  Multiplication combines two independent
errors, so their ratios multiply and the level must absorb a squaring --- the
same single-level cost, for the same reason, as composing two steps through a
midpoint.  Inversion combines nothing; it is a bijection of the carrier that
preserves multiplicative distance, so it is free.  Every level bookkeeping in
the rest of the file traces back to this one distinction: the constructions
that combine two streams take a $\ctor{join}$ of anchors and ask their
hypotheses for $\ctor{succ}\,\ell$, and the constructions that transform one
stream keep both the anchor and the level.
\end{remark}

\subsection{Operations on streams}

\begin{definition}[\lean{MulCauchyStream.mul} and \lean{MulCauchyStream.inv}]
\label{Algebra:def:stream-ops}
The product of two Cauchy streams is the pointwise product, and the inverse is
the pointwise inverse; each carries a proof that the result is again Cauchy.
The two $\ctor{simp}$ lemmas $\name{mul\_term}$ and $\name{inv\_term}$ record
that the terms are what they look like, both by $\ctor{rfl}$.
\end{definition}

\begin{proof}[Proof that the product is Cauchy]
Given a level $\ell$, ask each factor's Cauchy property for
$\ctor{succ}\,\ell$, obtaining anchors $a_a$ and $a_b$, and take
$\name{join}(a_a, a_b)$.  For $m$ and $n$ in that tail, the two bounds on
$\name{join}$ and transitivity of $\name{AtOrAfter}$, all from
\texttt{PrimeTensor/Completion.lean}, place both indices in both original
tails, so both factors are close at $\ctor{succ}\,\ell$, and
Theorem~\ref{Algebra:thm:scalewithin-mul} combines them at $\ell$.
\end{proof}

\begin{proof}[Proof that the inverse is Cauchy]
Ask the Cauchy property for $\ell$ itself, keep the same anchor, and apply
Theorem~\ref{Algebra:thm:scalewithin-inv} termwise.  No $\name{join}$ and no
$\ctor{succ}$ occur.
\end{proof}

\subsection{Compatibility with asymptotic equivalence}

\begin{lemma}[\lean{MulAsymptotic.of\_pointwise}]\label{Algebra:lem:of-pointwise}
If $a_n = b_n$ for every $n$, then $a$ and $b$ are asymptotic.
\end{lemma}

\begin{leancode}
theorem of_pointwise {a b : MulCauchyStream}
    (h : ∀ n : Depth, a.term n = b.term n) :
    MulAsymptotic a b := by
  intro level
  refine ⟨.one, ?_⟩
  intro n hn
  rw [h n]
  exact MulRat.scaleWithin_refl level (b.term n)
\end{leancode}

\begin{proof}
Take the anchor $\ctor{one}$, so the tail is everything, rewrite by the
pointwise equality, and apply unconditional reflexivity of
$\name{ScaleWithin}$.
\end{proof}

This small lemma is what makes stratum (iii) cheap.  Every group law below is
an equality that already holds term by term, so none of them needs a level, an
anchor or an estimate: $\name{of\_pointwise}$ turns the termwise identity
straight into the required asymptotic equivalence, and $\ctor{Quotient.sound}$
turns that into an equality in $\name{MulReal}$.

\begin{lemma}[\lean{MulAsymptotic.mul} and \lean{MulAsymptotic.inv}]
\label{Algebra:lem:congruence}
$\name{MulAsymptotic}$ is a congruence for both operations: from
$a \sim a'$ and $b \sim b'$ follows $ab \sim a'b'$, and from $a \sim b$ follows
$a^{-1} \sim b^{-1}$.
\end{lemma}

\begin{proof}
For the product, the same three moves as the transitivity proof of
\texttt{PrimeTensor/Completion.lean}: ask both hypotheses for
$\ctor{succ}\,\ell$, take $\name{join}$ of the anchors, and combine with
Theorem~\ref{Algebra:thm:scalewithin-mul}.  For the inverse, keep the anchor
and the level and apply Theorem~\ref{Algebra:thm:scalewithin-inv} termwise.
\end{proof}

\subsection{The operations on \texorpdfstring{$\name{MulReal}$}{MulReal}}

\begin{definition}[\lean{MulReal.mul} and \lean{MulReal.inv}]
\label{Algebra:def:mulreal-ops}
$\name{mul}$ is $\ctor{Quotient.liftOn_2}$ of the stream product and
$\name{inv}$ is $\ctor{Quotient.liftOn}$ of the stream inverse, each with
Lemma~\ref{Algebra:lem:congruence} as its well-definedness obligation.  They
are registered as the $\ctor{Mul}$ and $\ctor{Inv}$ instances of
$\name{MulReal}$.
\end{definition}

\begin{leancode}
def inv (a : MulReal) : MulReal :=
  Quotient.liftOn a
    (fun x => ofStream (MulCauchyStream.inv x))
    (by
      intro a b h
      change MulAsymptotic a b at h
      change
        Quotient.mk MulAsymptotic.setoid (MulCauchyStream.inv a) =
          Quotient.mk MulAsymptotic.setoid (MulCauchyStream.inv b)
      apply Quotient.sound
      exact MulAsymptotic.inv h)
\end{leancode}

\begin{theorem}[The group laws]\label{Algebra:thm:group}
$\name{MulReal}$ is a commutative group under $\name{mul}$ with unit
$1 = \name{ofRat}(1)$ and inverse $\name{inv}$.  Explicitly:
$1 \cdot a = a$ and $a \cdot 1 = a$; $(ab)c = a(bc)$; $ab = ba$;
$(a^{-1})^{-1} = a$; $a \cdot a^{-1} = 1$ and $a^{-1} \cdot a = 1$.
\end{theorem}

\begin{proof}
All but the last have the same four-line shape: induct to replace each element
by a representative stream, apply $\ctor{Quotient.sound}$ to reduce the claim
to an asymptotic equivalence, apply
Lemma~\ref{Algebra:lem:of-pointwise} to reduce that to a termwise equality, and
close the termwise equality with the corresponding law of
\texttt{PrimeTensor/MulRat.lean} --- $\name{one\_mul}$, $\name{mul\_one}$,
$\name{mul\_assoc}$, $\name{mul\_comm}$, $\name{inv\_inv}$, $\name{mul\_inv}$
respectively.  The unit cases additionally use that the constant stream at $1$
has every term $1$.  The last law, $a^{-1}a = 1$, is obtained from the previous
two by commuting, exactly as in \texttt{PrimeTensor/Scale/Laws.lean}.
\end{proof}

\begin{remark}[Where the analysis went]\label{Algebra:rem:no-analysis}
Not one of the seven laws mentions a level, an anchor, or
$\name{ScaleWithin}$.  That is not an accident of presentation: the laws hold
strictly, term by term, because they hold in $\name{MulRat}$, and equality of
streams is finer than asymptotic equivalence.  All of the analytic content of
the file was spent in
Theorems~\ref{Algebra:thm:scalewithin-mul} and~\ref{Algebra:thm:scalewithin-inv},
which are what make the operations well defined on the quotient at all; once
they are proved, the group structure descends for free.

The same will not be true of an order or of a limit, where the statement itself
mentions the scale and cannot be checked termwise.
\end{remark}

\begin{remark}[Still a bare $\ctor{Mul}$, not a $\ctor{CommGroup}$]
\label{Algebra:rem:no-instance}
The seven laws amount to $\name{MulReal}$ being an abelian group, but no
$\ctor{CommGroup}$ instance is declared: only $\ctor{One}$, $\ctor{Mul}$ and
$\ctor{Inv}$.  This repeats the choice made for $\name{MulRat}$, and it has the
same consequence, visible already in this file --- the regrouping of four
factors in Lemma~\ref{Algebra:lem:four-shuffle} has to be written by hand
rather than discharged by $\ctor{ac\_rfl}$ or $\ctor{group}$.  It keeps the
development independent of the library's algebraic hierarchy, at the price of
re-proving inside it whatever the hierarchy would have supplied.
\end{remark}

\begin{remark}[What is still absent]\label{Algebra:rem:absent}
There is no order on $\name{MulReal}$, no limit notion, and no statement that
$\name{MulReal}$ is complete for the scale it was built from --- the standard
next obligation for any completion.  Nor is $\name{ofRat}$ shown to be a
homomorphism or shown injective, although the first would now follow from
Lemma~\ref{Algebra:lem:of-pointwise} in one line, since the constant stream of
a product is the product of the constant streams.  Consistently with the
module header, nothing additive appears: no addition, no zero, no subtraction,
and no coefficient drawn from $\mathbb{Q}$ or $\mathbb{R}$.
\end{remark}

\subsection*{Summary of the correspondence}

\begin{center}
\small
\begin{tabular}{@{}lll@{}}
\hline
Lean declaration & Kind & Here \\
\hline
\lean{MulRat.inv\_mul\_pair}          & theorem     & Lem.~\ref{Algebra:lem:inv-mul-pair} \\
\lean{MulRat.inv\_one}                & simp thm    & Lem.~\ref{Algebra:lem:inv-mul-pair} \\
\lean{MulRat.mul\_four\_shuffle}      & private thm & Lem.~\ref{Algebra:lem:four-shuffle} \\
\lean{MulRat.ratio\_mul\_pair}        & theorem     & Lem.~\ref{Algebra:lem:ratio-mul} \\
\lean{MulRat.ratio\_inv\_pair}        & theorem     & Lem.~\ref{Algebra:lem:ratio-inv} \\
\lean{MulRat.scaleWithin\_mul}        & theorem     & Thm.~\ref{Algebra:thm:scalewithin-mul} \\
\lean{MulRat.scaleWithin\_inv}        & theorem     & Thm.~\ref{Algebra:thm:scalewithin-inv} \\
\lean{MulCauchyStream.mul}            & definition  & Def.~\ref{Algebra:def:stream-ops} \\
\lean{MulCauchyStream.inv}            & definition  & Def.~\ref{Algebra:def:stream-ops} \\
\lean{MulCauchyStream.mul\_term}      & simp thm    & Def.~\ref{Algebra:def:stream-ops} \\
\lean{MulCauchyStream.inv\_term}      & simp thm    & Def.~\ref{Algebra:def:stream-ops} \\
\lean{MulAsymptotic.of\_pointwise}    & theorem     & Lem.~\ref{Algebra:lem:of-pointwise} \\
\lean{MulAsymptotic.mul}              & theorem     & Lem.~\ref{Algebra:lem:congruence} \\
\lean{MulAsymptotic.inv}              & theorem     & Lem.~\ref{Algebra:lem:congruence} \\
\lean{MulReal.mul}                    & definition  & Def.~\ref{Algebra:def:mulreal-ops} \\
\lean{MulReal.inv}                    & definition  & Def.~\ref{Algebra:def:mulreal-ops} \\
\lean{Mul MulReal}                    & instance    & Def.~\ref{Algebra:def:mulreal-ops} \\
\lean{Inv MulReal}                    & instance    & Def.~\ref{Algebra:def:mulreal-ops} \\
\lean{MulReal.one\_mul}               & simp thm    & Thm.~\ref{Algebra:thm:group} \\
\lean{MulReal.mul\_one}               & simp thm    & Thm.~\ref{Algebra:thm:group} \\
\lean{MulReal.mul\_assoc}             & theorem     & Thm.~\ref{Algebra:thm:group} \\
\lean{MulReal.mul\_comm}              & theorem     & Thm.~\ref{Algebra:thm:group} \\
\lean{MulReal.inv\_inv}               & simp thm    & Thm.~\ref{Algebra:thm:group} \\
\lean{MulReal.mul\_inv}               & simp thm    & Thm.~\ref{Algebra:thm:group} \\
\lean{MulReal.inv\_mul}               & simp thm    & Thm.~\ref{Algebra:thm:group} \\
\hline
\end{tabular}
\end{center}

\noindent
Every declaration of \texttt{PrimeTensor/Algebra.lean} appears above.  The file
contains no \lean{sorry}, no axiom, and no unproved named proposition.
Remarks~\ref{Algebra:rem:level-cost}, \ref{Algebra:rem:no-analysis},
\ref{Algebra:rem:no-instance} and~\ref{Algebra:rem:absent} are stated for the
reader and are not declarations of the file.
